\chapter{Literature Review}

\section{Bayesian Logistic Regression for Binary Response Data}

Logistic regression is a standard member of the generalized linear model family for binary and binomial response data. In the generalized linear model framework, a response distribution from the exponential family is connected to a linear predictor through a link function; for binary response modelling, the logit link is commonly used to model the log-odds of the event probability as a linear function of covariates \cite{mccullagh1989generalized,hosmer2013applied}. This makes logistic regression suitable for modelling whether an observation belongs to an event class or a reference class.

For a binary response $y_i\in\{0,1\}$, logistic regression assumes
\[
y_i\mid p_i \sim \operatorname{Bernoulli}(p_i),
\qquad
\operatorname{logit}(p_i)
=
\log\left(\frac{p_i}{1-p_i}\right)
=
\mathbf{x}_i^T\boldsymbol{\beta}.
\]
The regression coefficient $\beta_j$ is interpreted on the log-odds scale, while $\exp(\beta_j)$ is interpreted as an odds ratio for a one-unit change in the corresponding predictor or relative to a reference category for categorical predictors \cite{hosmer2013applied}. When categorical predictors are included, one category is treated as the reference category, and the coefficients for the remaining categories describe log-odds contrasts relative to that baseline.

Bayesian logistic regression extends this model by treating the coefficient vector $\boldsymbol{\beta}$ as unknown and assigning it a prior distribution. The posterior distribution is then proportional to the product of the likelihood and the prior distribution,
\[
\pi(\boldsymbol{\beta}\mid \mathbf{y},\mathbf{X})
\propto
L(\mathbf{y}\mid \boldsymbol{\beta},\mathbf{X})\pi(\boldsymbol{\beta}).
\]
This posterior distribution represents uncertainty about the regression coefficients after observing the data \cite{gelman2013bayesian}. In the present study, Bayesian logistic regression provides the posterior target that is later approximated using full-data Metropolis--Hastings sampling and Consensus Monte Carlo.

\section{Prior Specification in Bayesian Logistic Regression}

Prior specification is an important part of Bayesian regression modelling. In logistic regression, coefficient estimates may become unstable when the data are sparse, when predictors are imbalanced, or when separation occurs. Weakly informative priors are designed to regularise estimation without imposing strong prior beliefs. Gelman et al. proposed weakly informative Student-$t$ prior distributions for logistic and related regression models and recommended a Cauchy prior with centre $0$ and scale $2.5$ as a default prior choice after suitable scaling of predictors \cite{gelman2008weakly}.

Following this motivation, the present study assigns independent Cauchy priors to the regression coefficients,
\[
\beta_j \sim \operatorname{Cauchy}(0,2.5),
\qquad j=0,\ldots,p.
\]
The prior is centred at zero, representing no strong prior direction for the coefficient, while its heavy tails allow large coefficient values when supported by the data. This makes the prior weakly informative rather than strongly restrictive. However, the heavy-tailed nature of the Cauchy prior also requires care. Ghosh, Li and Mitra examined theoretical and computational issues associated with Cauchy priors in Bayesian logistic regression, particularly under separation, and showed that posterior behaviour may be problematic in some cases \cite{ghosh2015cauchy}. Thus, the Cauchy prior is useful as a weakly informative default, but its computational implications should be acknowledged when interpreting posterior simulation output.

\section{Markov Chain Monte Carlo and Metropolis--Hastings}

For many Bayesian models, posterior summaries cannot be obtained analytically because the posterior normalising constant is unavailable or posterior expectations require high-dimensional integration. Markov Chain Monte Carlo methods address this problem by constructing a Markov chain whose stationary distribution is the target posterior distribution. After burn-in, simulated draws from the chain can be used to approximate posterior means, posterior standard deviations, quantiles, credible intervals, and other summaries. Tierney provided a general theoretical treatment of using Markov chains to explore posterior distributions \cite{tierney1994markov}.

The Metropolis algorithm was introduced by Metropolis et al. for simulation-based computation \cite{metropolis1953equation}, and Hastings later generalized it to allow more general proposal distributions \cite{hastings1970monte}. The resulting Metropolis--Hastings algorithm is a flexible MCMC method for sampling from a target distribution known up to a proportionality constant. Chib and Greenberg provide a clear exposition of the algorithm and its practical implementation \cite{chib1995understanding}.

At iteration $t$, a candidate value $\boldsymbol{\beta}^{*}$ is generated from a proposal density $q(\boldsymbol{\beta}^{*}\mid\boldsymbol{\beta}^{(t-1)})$. The candidate is accepted with probability
\[
\alpha
=
\min\left\{
1,
\frac{
\pi(\boldsymbol{\beta}^{*}\mid \mathbf{y},\mathbf{X})
q(\boldsymbol{\beta}^{(t-1)}\mid\boldsymbol{\beta}^{*})
}{
\pi(\boldsymbol{\beta}^{(t-1)}\mid \mathbf{y},\mathbf{X})
q(\boldsymbol{\beta}^{*}\mid\boldsymbol{\beta}^{(t-1)})
}
\right\}.
\]
This general form covers both Random Walk Metropolis--Hastings and Independent Metropolis--Hastings. In Random Walk Metropolis--Hastings, the proposal is centred at the current state, so the proposal is symmetric when a multivariate normal random-walk proposal is used. In Independent Metropolis--Hastings, candidates are generated from a fixed proposal distribution independent of the current state, so the proposal-density correction must be retained in the acceptance probability.

\section{Proposal Design and MCMC Accuracy}

The efficiency of Metropolis--Hastings sampling depends strongly on the proposal distribution. If the proposal variance is too small, the chain may have a high acceptance rate but move slowly through the posterior distribution. If the proposal variance is too large, many proposed values may be rejected. Therefore, proposal scaling is a central practical issue in Random Walk Metropolis--Hastings.

For multivariate random-walk proposals, Roberts, Gelman and Gilks developed optimal-scaling results that motivate proposal covariance matrices proportional to
\[
\frac{2.38^2}{d}\widehat{\boldsymbol{\Sigma}},
\]
where $d$ is the parameter dimension and $\widehat{\boldsymbol{\Sigma}}$ is an estimate of the target covariance structure \cite{roberts1997weak}. In the present implementation, this scaling is combined with a tuning factor and an estimated covariance matrix from a preliminary frequentist logistic regression fit. Cholesky decomposition is used as a computational device for generating correlated multivariate normal proposal draws. If $\boldsymbol{\Sigma}=\mathbf{L}\mathbf{L}^{T}$, where $\mathbf{L}$ is a Cholesky factor, then correlated normal proposals can be generated by transforming independent standard normal random variables using $\mathbf{L}$.

The Independent Metropolis--Hastings sampler uses a fixed multivariate normal proposal centred at the frequentist logistic regression estimate. Its proposal covariance is also based on the estimated covariance matrix from the frequentist fit, multiplied by a tuning factor. This design uses the frequentist approximation to provide a proposal distribution with a plausible centre and scale, while the Metropolis--Hastings acceptance probability corrects for the difference between the proposal distribution and the target posterior distribution.

Monte Carlo accuracy must also be assessed because MCMC samples are autocorrelated. Effective sample size measures the amount of independent-sample information represented by an autocorrelated chain. Monte Carlo standard error measures the simulation uncertainty of posterior estimates. Flegal, Haran and Jones emphasized that reporting Monte Carlo standard errors is important for assessing the numerical reliability of MCMC estimates \cite{flegal2008mcmc}. In this study, ESS and MCSE are reported together with posterior summaries to evaluate the accuracy of the full-data and consensus posterior estimates.

\section{Multi-chain Convergence Diagnostics}

Running multiple independent chains provides additional information about whether MCMC simulations have reached similar stationary behaviour. Gelman and Rubin proposed a multiple-sequence convergence diagnostic based on comparing between-chain and within-chain variation \cite{gelman1992inference}. This diagnostic is commonly expressed as the potential scale reduction factor, or $\hat{R}$, where values close to $1$ indicate that the independent chains have similar behaviour. Brooks and Gelman extended the diagnostic framework and discussed general methods for monitoring convergence in iterative simulations \cite{brooks1998general}.

The multi-chain implementation in this project reports Gelman--Rubin diagnostics for full-data chains, shard-level chains, and final CMC consensus chains. It also computes ESS using the chains as a \texttt{coda::mcmc.list} object rather than computing autocorrelation-based diagnostics from artificially pooled chains. This distinction is important because pooling several chains into one long matrix may be appropriate for posterior summaries, but it should not be treated as a single genuine time series for convergence diagnostics.

More recent work by Vehtari et al. proposed rank-normalized, folded, and localized versions of $\hat{R}$, showing that traditional $\hat{R}$ can miss convergence problems in heavy-tailed or scale-varying chains \cite{vehtari2021rank}. The present study does not implement these newer rank-normalized diagnostics; instead, it uses the classical Gelman--Rubin diagnostic through the \texttt{coda} package. Nevertheless, the modern diagnostic literature is useful for interpreting the limitations of $\hat{R}$ and for motivating the combined reporting of $\hat{R}$, ESS, MCSE, trace plots, and autocorrelation plots.

\section{Distributed Bayesian Computation}

Full-data MCMC can become computationally expensive for large datasets because each iteration may require evaluating the likelihood over all observations. Distributed Bayesian computation addresses this issue by partitioning the data, performing computation separately on subsets, and recombining the subset-level results. Huang and Gelman proposed an early subset-based approach for Bayesian computation with large datasets, using separate analyses on data partitions followed by approximate combination of the resulting inferences \cite{huang2005sampling}. This work is relevant because it identifies the main statistical challenge that later Consensus Monte Carlo methods address: local computations are easy to parallelise, but the local posterior information must be recombined in a way that approximates the full posterior.

In the present study, data partitioning is not used merely for data management. It defines the computational structure of the Consensus Monte Carlo method. The observations are divided into $K=10$ shards using stratified assignment by response class so that each shard retains representation from both binary response categories. Shard-level MCMC computations are then run independently, allowing the expensive likelihood evaluations to be distributed across workers.

\section{Consensus Monte Carlo}

Consensus Monte Carlo is a distributed Bayesian computation method designed to approximate a full-data posterior distribution using subset posterior simulations. Scott et al. proposed the Consensus Monte Carlo algorithm as a low-communication approach for Bayesian inference with large datasets \cite{scott2016bayes}. The key idea is to divide the data into shards, run separate Monte Carlo simulations on each shard, and then combine the resulting subposterior draws to approximate the full posterior distribution.

Suppose the full likelihood factorises over $K$ data shards as
\[
L(\mathbf{y}\mid\boldsymbol{\beta},\mathbf{X})
=
\prod_{k=1}^{K}
L(\mathbf{y}_k\mid\boldsymbol{\beta},\mathbf{X}_k).
\]
A common Consensus Monte Carlo construction defines the $k$th subposterior as
\[
\pi_k(\boldsymbol{\beta}\mid\mathbf{y}_k,\mathbf{X}_k)
\propto
L(\mathbf{y}_k\mid\boldsymbol{\beta},\mathbf{X}_k)
\pi(\boldsymbol{\beta})^{1/K}.
\]
The fractional prior contribution ensures that multiplying all subposterior kernels recovers a kernel proportional to the full posterior. This construction is used in the present methodology for both CMC-RWMH and CMC-IMH.

The remaining problem is how to combine subposterior samples. In approximately Gaussian settings, precision-weighted averaging is a natural strategy. If $\widehat{\boldsymbol{\Sigma}}_k$ is the estimated covariance matrix of subposterior $k$ and $\mathbf{W}_k=\widehat{\boldsymbol{\Sigma}}_k^{-1}$ is the corresponding precision matrix, then consensus draws can be formed using precision-weighted combinations of shard-level draws. Scott compared several Consensus Monte Carlo strategies and emphasized that the practical performance of CMC depends not only on generating local posterior samples but also on the aggregation rule used to combine them \cite{scott2017comparing}. This is directly relevant to the present study, which uses precision-weighted combination of subposterior samples.

\section{Evaluation of Consensus Monte Carlo Approximation}

Because Consensus Monte Carlo is an approximation to full-data posterior inference, its output must be compared with a full-data reference. In this study, full-data Random Walk Metropolis--Hastings serves as the reference for CMC-RWMH, and full-data Independent Metropolis--Hastings serves as the reference for CMC-IMH. Posterior means, medians, posterior standard deviations, credible intervals, odds ratios, acceptance rates, $\hat{R}$, ESS, MCSE, trace plots, autocorrelation plots, density plots, and runtime are used to compare the methods.

The posterior mean root mean squared error is used as a direct numerical measure of approximation error. For a CMC posterior mean vector $\widehat{\boldsymbol{\mu}}_{\mathrm{CMC}}$ and a full-data posterior mean vector $\widehat{\boldsymbol{\mu}}_{\mathrm{full}}$, the RMSE is
\[
\operatorname{RMSE}
=
\sqrt{
\frac{1}{d}
\sum_{j=1}^{d}
\left(
\widehat{\mu}_{\mathrm{CMC},j}
-
\widehat{\mu}_{\mathrm{full},j}
\right)^2
}.
\]
A smaller RMSE indicates that the consensus posterior mean is closer to the corresponding full-data posterior mean. However, RMSE alone is not sufficient. It must be interpreted together with diagnostic measures and computational runtime because a method may be fast but inaccurate, or accurate but computationally inefficient. Therefore, the evaluation of CMC in this study considers both statistical agreement with the full-data posterior and computational efficiency.
